% Figure legends (Task 46): moved out of the figure floats, listed after Additional information as the journal asks.
% At most 350 words each. Every number is a macro.
\section*{Figure legends}
\label{sec:legends}

\noindent\textbf{Figure~\mainFigPipeline.}
Study design.
FT-Transformer classifiers (\nSeeds{} seeds) and an XGBoost reference were trained on LYCOS-IDS2017 and
LycoS-Unicas-IDS2018.
Clean explanations compare three attention variants with IG, GradientSHAP, occlusion and TreeSHAP by agreement and by
deletion and insertion faithfulness.
Constrained attacks (A0 random control, A1 evasion, A1b targeted, A2 attention, A3 evasion and attention) change only
the controls of a flow within their budgets, recompute the derived features and are checked by a separate validity
checker; their effect on attention and on faithfulness is measured on the attacked rows.
The detectors D (confidence and distance scores) and D$_+$ (D plus \nConsistency{} attention-attribution consistency
scores) and \nControlsAll{} control detectors are evaluated against the non-adaptive attacker A6 and against adaptive
attackers: A8 (differentiable stand-ins for the consistency scores), A8x (the exact score of D$_+$, with surrogate
gradients only for the \nRankScores{} rank-based scores), A6x (A6 with the iterate selection of A8x) and an attacker of
each control that models its own scores, compared under the equal-effort union.

\medskip
\noindent\textbf{Figure~\mainFigHeadline.}
D$_+$ minus D in the share of successful and undetected source rows under the equal-effort union at
$\epsilon = \paramEpsHead$ on the primary row set V1, per dataset.
A positive value means that D$_+$ left more attack rows successful and undetected than D.
Top row: D$_+$ faces A1 and the three settings of the stand-in attacker A8, D faces A1 and A6.
Middle row: D$_+$ faces A1 and the three settings of the exact-objective attacker A8x, D faces A1 and A6.
Bottom row: as the middle row, with D facing A1 and A6x, which gives D's attacker the iterate selection of A8x.
Filled marks: difference pooled over the \nSeeds{} seeds with the \paramCILevel{} paired percentile bootstrap interval
over source rows within seed (\paramBoot{} draws); the printed values are the pooled difference and its interval.
Open marks: the differences of seeds 0, 1 and 2 (labels s0, s1, s2).
The intervals do not capture variation between seeds; \axHeadOppSeedsPhrase{} had the opposite sign in the middle and bottom rows.
A8x uses surrogate gradients for the \nRankScores{} rank-based consistency scores.
On LycoS-Unicas-IDS2018, D alone left only \axHeadVoneEighteenD{} of the source rows successful and undetected, so the
difference there sits near a floor.
Source rows per seed: \nSrcVoneSeventeen{} on LYCOS-IDS2017 and \nSrcVoneEighteen{} on LycoS-Unicas-IDS2018.
